% ============================================================
\section{Iteración 4 --- DRV-03: presupuesto de un segundo por jugada}\label{sec:it04}
% ============================================================

% ------------------------------------------------------------
\subsection{Driver y alcance}

\begin{tabla}{|L{3.0cm}|Y|}
\fichacab
\parte{Driver}{DRV-03 --- Presupuesto de un segundo por jugada.}
\parte{Enunciado}{Cada jugada en línea dispone de menos de un segundo, con mediana bajo 300 ms, repartido de forma explícita entre validación, red y persistencia. Nada entra en el camino síncrono de la respuesta sin presupuesto asignado.}
\parte{Por qué le toca este turno}{Suma 15 de 20 y agrupa ASR-01, el único requisito con un número impuesto. Es el driver que impide maximizar los demás, así que solo se puede repartir cuando existe lo que consume el tiempo: el arbitraje de la iteración~1, el canal protegido de la iteración~2 y el motor de la iteración~3. Antes de eso, cualquier reparto sería una suposición.}
\parte{Qué decide esta iteración}{Cuánto tiempo puede consumir cada tramo del camino de una jugada y qué se hace para que la suma quede por debajo del segundo, con mediana bajo 300 ms.}
\parte{Decisiones ya tomadas}{DEC-01 a DEC-03 y lo decidido en las tres iteraciones anteriores. Pesa especialmente que DEC-02 haya descartado la validación en el cliente: cada jugada cuesta un viaje completo de red, y el margen que queda para el servidor es lo que esta iteración tiene que repartir.}
\parte{Qué no decide}{Qué se persiste y cuándo, que es de la iteración~5, aunque aquí se fije cuánto tiempo puede costar. Tampoco vuelve a discutir la autoridad, el protocolo ni las reglas.}
\end{tabla}

% ------------------------------------------------------------
\subsection{Paso 1. Revisión de entradas}

\subsubsection*{ASR que agrupa el driver}

Del conjunto de la sección~\ref{sec:asr}, estos son los ASR que DRV-03 agrupa y el papel que cumple cada uno en esta iteración.

\begin{tabla}{|L{1.4cm}|Y|L{1.7cm}|L{3.6cm}|}
\hline
\filacab \cab{ASR} & \cab{Qué debe cumplir la arquitectura} & \cab{Prioridad} & \cab{Papel en esta iteración} \\ \hline
\endhead
ASR-01 & Una jugada en línea se responde en menos de 1 segundo, con mediana bajo 300 ms, con el servidor decidiendo y el tráfico cifrado. & Prioritario & Objetivo. Su medida es ESC-01. \\ \hline
ASR-09 & Una partida en curso se puede observar, incluso sin cuenta, recibiendo el estado con retraso y sin frenar a los jugadores. & Media & Condición: la notificación a los espectadores no puede competir con la de los participantes. \\ \hline
\end{tabla}

\subsubsection*{Objetivos de diseño}

\begin{tabla}{|L{1.5cm}|L{5.0cm}|Y|}
\hline
\filacab \cab{ID} & \cab{Objetivo} & \cab{Origen y qué exige aquí} \\ \hline
\endhead
OBJ-12 & El jugador recibe respuesta inmediata a su jugada. &
CRN-02 y REQ-02, con STK-01. Es el único objetivo con un número impuesto de antemano. \\ \hline
OBJ-13 & La partida se siente igual de fluida con el servidor decidiendo y el tráfico protegido. &
TEN-01. Lo que se ganó en integridad y confidencialidad no puede pagarse con una partida lenta. \\ \hline
OBJ-14 & Lo que el juego registra para poder auditarse no retrasa al jugador. &
TEN-06, con STK-13 y STK-16. \\ \hline
\end{tabla}

\subsubsection*{Requisitos funcionales primarios}

\begin{tabla}{|L{1.3cm}|L{4.2cm}|Y|}
\hline
\filacab \cab{CU} & \cab{Caso de uso} & \cab{Qué exige del camino de la jugada} \\ \hline
\endhead
CU-21 & Colocar un muro &
Es el peor caso: el servidor comprueba solape y camino a la meta de todos los jugadores sobre el tablero vigente, con el tablero avanzado y la mayoría de los muros ya colocados. \\ \hline
CU-20 & Mover el peón &
Es el caso frecuente y más barato, pero recorre el mismo camino completo y compite por el mismo segundo. \\ \hline
CU-34 & Ver una partida como espectador &
Recibe el estado con un retraso previsto, así que su notificación no puede competir con la de los participantes. \\ \hline
CU-28 & Enviar un mensaje de chat &
Viaja por el mismo canal y el mismo servicio de protocolo durante la partida. \\ \hline
\end{tabla}

\subsubsection*{Escenarios de atributos de calidad}

\begin{tabla}{|L{1.7cm}|L{3.3cm}|Y|}
\hline
\filacab \cab{Escenario} & \cab{Atributo} & \cab{Medida} \\ \hline
\endhead
ESC-01 \newline (ASR-01) & Desempeño: comportamiento temporal &
Con 16 muros colocados, 150 ms de ida y vuelta y 20 partidas simultáneas: el 100~\% de 1\,000 colocaciones responde en menos de 1 segundo, con mediana bajo 300 ms; el oponente la ve en el mismo plazo; 0 jugadas aceptadas se revierten. \\ \hline
ESC-05, ESC-06 & Seguridad &
Condición. Las medidas de las iteraciones~1 y~2 se vuelven a ejecutar con la protección activa: ninguna optimización puede lograrse quitando comprobaciones. \\ \hline
ESC-09 & Fiabilidad &
Condición. La ventana de cinco segundos de escritura diferida es la que hace posible el presupuesto; si la iteración~5 la reduce, el reparto cambia. \\ \hline
\end{tabla}

\subsubsection*{Restricciones}

\begin{tabla}{|L{1.9cm}|L{4.3cm}|Y|}
\hline
\filacab \cab{ID} & \cab{Restricción} & \cab{Cómo limita esta iteración} \\ \hline
\endhead
REQ-02 & El tiempo de respuesta debe ser menor a 1 segundo & Es el número impuesto que hay que repartir. \\ \hline
REQ-01 & Modo multijugador en línea obligatorio & La red está siempre en el camino: no hay modo local que la evite. \\ \hline
CON-02 & La comunicación es por TCP & La entrega ordenada y confiable cuesta retransmisiones cuando la red pierde paquetes. \\ \hline
DEP-04 & Infraestructura externa & La latencia entre jugador y servidor no la controla el equipo; el ambiente de la medida la fija en 150 ms. \\ \hline
CON-13 & 14 semanas & Descarta optimizaciones caras de construir: lo que no quepa en el plazo no sirve. \\ \hline
\end{tabla}

\subsubsection*{Decisiones de proyecto ya tomadas}

\begin{tabla}{|L{1.5cm}|L{4.3cm}|Y|}
\hline
\filacab \cab{ID} & \cab{Decisión} & \cab{Qué impone a esta iteración} \\ \hline
\endhead
DEC-02 & Toda la lógica del juego en el servidor & Descarta la validación previa en el cliente, que es la forma habitual de bajar la mediana. El viaje completo de red es obligatorio. \\ \hline
CD-12 & Canal seguro sobre el socket & Agrega un costo por mensaje que hay que medir y presupuestar. \\ \hline
CD-10 & Escritura fuera del camino de la respuesta & Saca la base de datos del tramo síncrono, que es lo que deja margen. \\ \hline
CD-21 & El motor enumera jugadas legales y explica el rechazo & El trabajo del motor por jugada es mayor que una simple validación, y el mensaje de turno es más grande. \\ \hline
\end{tabla}

\subsubsection*{Concerns}

\begin{tabla}{|L{1.5cm}|Y|}
\hline
\filacab \cab{ID} & \cab{Concern y qué aporta a la iteración} \\ \hline
\endhead
CRN-02 & \emph{Cuando coloco una barrera quiero que el juego me conteste al instante. Si tarda más de un segundo en decirme si vale, la partida se me hace pesada.} \\ \hline
CRN-10, CRN-22 & Recuerdan que el servidor autoritativo y el cifrado son lo que consume el presupuesto, y que no se pueden retirar para ganar tiempo. \\ \hline
CRN-19 & La persistencia compite por el mismo segundo; su política se decide después, pero su costo se presupuesta aquí. \\ \hline
\end{tabla}

\subsubsection*{Preguntas abiertas que condicionan la iteración}

\begin{tabla}{|L{1.3cm}|Y|}
\hline
\filacab \cab{ID} & \cab{Cómo condiciona} \\ \hline
\endhead
Q-01 & Si habrá IA, su tiempo de cálculo tendrá que caber en el mismo segundo, porque ocupa el asiento de un jugador. \\ \hline
Q-02 & En LAN la latencia sería mucho menor, pero el presupuesto se fija para el caso en línea, que es el obligatorio. \\ \hline
\end{tabla}

% ------------------------------------------------------------
\subsection{Paso 2. Objetivo de la iteración y selección de entradas}

\subsubsection*{Priorización de las entradas}

\begin{tabla}{|L{3.2cm}|L{1.3cm}|Y|L{2.2cm}|}
\hline
\filacab \cab{Entrada} & \cab{Peso} & \cab{Por qué pesa lo que pesa} & \cab{Decisión} \\ \hline
\endhead
ESC-01 (ASR-01) & Alto & Es la medida del objetivo, con el peor caso ya definido: 16 muros, 150 ms y 20 partidas. & Se atiende \\ \hline
CU-21 & Alto & Es la jugada que más trabajo exige y la que la medida usa. & Se atiende \\ \hline
DEC-02 & Alto & Determina que el viaje de red completo es obligatorio, que es la mayor partida del presupuesto. & Se atiende \\ \hline
CD-10, CD-12, CD-21 & Alto & Son los tres costos ya decididos que hay que acomodar: cifrado, mensaje de turno grande y escritura diferida. & Se atiende \\ \hline
REQ-02, DEP-04 & Alto & Fijan el número y la latencia del ambiente. & Se atiende \\ \hline
ESC-05, ESC-06 & Medio & Condiciones: se vuelven a medir para comprobar que no se ganó tiempo quitando protección. & Se atiende como condición \\ \hline
CU-20, CU-34, CU-28 & Medio & Comparten el camino y el canal, y compiten por el mismo servidor. & Se atiende \\ \hline
ESC-09 & Medio & La ventana de cinco segundos es lo que permite el reparto; su definición es de la iteración~5. & Se atiende como supuesto \\ \hline
Q-01 & Bajo & Si la IA entra, consumirá parte del mismo presupuesto, pero hoy no existe. & Se atiende como condición \\ \hline
CON-13 & Bajo & Descarta optimizaciones que no quepan en el plazo. & Se atiende como filtro \\ \hline
\end{tabla}

\subsubsection*{Objetivo de la iteración}

\begin{tabla}{|L{3.0cm}|Y|}
\fichacab
\parte{Objetivo}{Que una jugada se responda en menos de un segundo, con mediana bajo 300 ms, con el servidor decidiendo y el tráfico protegido.}
\parte{ASR que satisface}{ASR-01, con ASR-09 como condición.}
\parte{Driver que cubre}{DRV-03.}
\parte{Medida}{La de ESC-01: con 16 muros colocados, 150 ms de ida y vuelta y 20 partidas simultáneas, el 100~\% de 1\,000 colocaciones responde en menos de 1 segundo, la mediana queda bajo 300 ms, el oponente ve la jugada en el mismo plazo y 0 jugadas aceptadas se revierten.}
\parte{Límite que respeta}{ESC-05 y ESC-06. Las medidas de seguridad de las iteraciones~1 y~2 se vuelven a ejecutar al final: el presupuesto no se cumple quitando comprobaciones.}
\parte{Incremento que produce}{La partida protegida de la iteración~2, ahora instrumentada: cada jugada deja registrado cuánto tardó en cada tramo, y una prueba automatizada con latencia simulada y 20 partidas muestra la distribución completa de los 1\,000 tiempos. El resultado es una tabla de presupuesto con lo medido frente a lo asignado.}
\parte{Criterio de éxito}{La iteración termina cuando la prueba cumple el máximo y la mediana, y las medidas de ESC-05 y ESC-06 siguen pasando.}
\end{tabla}

% ------------------------------------------------------------
\subsection{Paso 3. Elemento a descomponer}

% ------------------------------------------------------------
\Needspace*{0.4\textheight}
\subsubsection*{EL-04 --- El camino de la jugada}

\begin{tabla}{|L{2.6cm}|Y|}
\fichacab
\parte{Elemento}{El camino de la jugada: la colaboración completa que va desde que el jugador confirma hasta que él y su oponente ven el resultado, atravesando el cliente, el servicio de protocolo, el servicio de estado, el motor y la cola de escritura.}
\parte{Qué abarca}{Los tramos que consumen tiempo: la preparación y el cifrado del mensaje en el cliente, el viaje de ida, la validación de formato, la autorización, la ejecución del motor con la comprobación de camino, la actualización del estado, la construcción del mensaje de turno con las jugadas legales, el viaje de vuelta y la notificación al oponente. No abarca la escritura en la base de datos, que CD-10 sacó del camino.}
\parte{Por qué se elige}{Porque la medida del objetivo no pertenece a ningún componente por separado: ninguno de ellos puede cumplir el segundo por su cuenta ni incumplirlo por su cuenta. Se elige frente a descomponer solo el servicio de estado, que es donde está el trabajo de cómputo pero no el de red, y frente al servicio de protocolo, que aporta el costo por mensaje pero no decide cuántos mensajes hay. Es una colaboración, y esa es exactamente la unidad en la que se reparte un presupuesto.}
\parte{Qué falta decidir sobre él}{Cuánto se le asigna a cada tramo y cómo se comprueba que cada uno se mantiene dentro. Cuántos mensajes hay por jugada, porque cada viaje cuesta la latencia completa. Y qué se hace con el trabajo que no cabe: sacarlo del camino, hacerlo antes o hacerlo en paralelo.}
\parte{Evidencia en los casos de uso}{CU-21 enumera veintidós pasos entre la confirmación del jugador y el dibujo del muro, y entre ellos están las dos comprobaciones caras, el solape y el camino a la meta de todos los jugadores (pasos 13 y 14), y la notificación a participantes y espectadores (paso 20). El caso de uso describe qué ocurre, pero no cuánto puede tardar nada de eso, que es lo que falta.}
\parte{Trazabilidad}{DRV-03. DEC-02. CD-05, CD-10, CD-12, CD-21. ESC-01 (ASR-01), con ESC-05, ESC-06 y ESC-09 como condiciones. CU-20, CU-21, CU-28, CU-34. REQ-02, REQ-01, CON-02, DEP-04. CRN-02, CRN-10, CRN-19. TEN-01, TEN-06. STK-01, STK-09, STK-05.}
\end{tabla}

% ------------------------------------------------------------
\subsection{Paso 4. Conceptos de diseño}

% ------------------------------------------------------------
\Needspace*{0.3\textheight}
\subsubsection*{CD-23 --- Presupuesto repartido por tramo, con tope asignado}

\begin{tabla}{|L{2.6cm}|Y|}
\fichacab
\parte{Estilo}{Táctica de desempeño: asignación explícita de recursos.}
\parte{Descripción}{El segundo se reparte de antemano entre los tramos del camino, y cada uno recibe un tope que no puede superar: la red se lleva la mayor parte por el ambiente de la medida, y lo que queda se distribuye entre el protocolo, el arbitraje, el motor y la notificación. Atiende ESC-01 al convertir un número global en obligaciones locales que se pueden comprobar por separado. Se elige frente a medir solo el total, que dice que se incumple pero no dónde.}
\parte{Efectos}{Cualquier trabajo nuevo tiene que pedir presupuesto a alguien, lo que convierte al segundo en un criterio de diseño y no en una comprobación final. Da a cada iteración posterior un límite concreto: la persistencia de la iteración~5 sabe cuánto puede costar antes de tener que salir del camino.}
\parte{Trazabilidad}{DRV-03. ESC-01 (ASR-01). CU-20, CU-21. REQ-02, DEP-04. CRN-02. TEN-01, TEN-06. STK-01, STK-05.}
\end{tabla}

% ------------------------------------------------------------
\Needspace*{0.3\textheight}
\subsubsection*{CD-24 --- Un solo viaje de ida y vuelta por jugada}

\begin{tabla}{|L{2.6cm}|Y|}
\fichacab
\parte{Estilo}{Táctica de desempeño: reducir la demanda sobre el recurso escaso.}
\parte{Descripción}{Una jugada se resuelve con un mensaje de ida y uno de vuelta, y todo lo que el cliente necesitará después viaja en ellos, incluidas las jugadas legales del siguiente turno (CD-21). Atiende ESC-01 en su parte más pesada: con 150 ms de ida y vuelta, cada viaje adicional consumiría la mitad del presupuesto de la mediana. Se elige frente a consultar al servidor mientras el jugador arrastra el muro por el tablero, que es lo que haría falta si el cliente no recibiera por adelantado lo que puede resaltar.}
\parte{Efectos}{El mensaje de turno es grande, lo que carga el costo sobre el protocolo en lugar de sobre la latencia, que es el intercambio favorable. Obliga a que el motor entregue en una sola llamada todo lo que el turno necesita. Fija un criterio para rechazar funcionalidades que exijan viajes extra dentro del turno.}
\parte{Trazabilidad}{DRV-03, DRV-02. DEC-02. CD-11, CD-13, CD-21. ESC-01. CU-20, CU-21 (pasos 3 a 7). REQ-02, CON-02. CRN-02, CRN-01. STK-01, STK-09.}
\end{tabla}

% ------------------------------------------------------------
\Needspace*{0.3\textheight}
\subsubsection*{CD-25 --- Cálculo incremental del camino a la meta}

\begin{tabla}{|L{2.6cm}|Y|}
\fichacab
\parte{Estilo}{Táctica de desempeño: reducir el trabajo de cómputo.}
\parte{Descripción}{La comprobación de que todos los jugadores conservan camino a su meta (RN-05 del CU-21) se apoya en el resultado de la jugada anterior en lugar de recorrer el tablero desde cero cada vez, y la enumeración de surcos válidos se calcula una vez por turno y se reutiliza para todas las posiciones candidatas. Atiende ESC-01 en el tramo del motor, que es el único trabajo de cómputo apreciable del camino. Se elige frente a recalcular todo en cada consulta, que multiplica el trabajo justo en el peor caso de la medida, con 16 muros colocados.}
\parte{Efectos}{El trabajo del motor por jugada queda acotado y medible, que es lo que CD-23 necesita para asignarle un tope. Obliga a que el motor mantenga una estructura derivada del tablero, que CD-20 exige que sea determinista y CD-05 mantiene dentro de la partida. Si Q-08 cambia las reglas, esa estructura tiene que revisarse.}
\parte{Trazabilidad}{DRV-03, DRV-05. CD-05, CD-20, CD-21. ESC-01, ESC-12. CU-21 (pasos 13 y 14, RN-04 y RN-05). REQ-02. CRN-02, CRN-13. Q-08. STK-01, STK-06.}
\end{tabla}

% ------------------------------------------------------------
\Needspace*{0.3\textheight}
\subsubsection*{CD-26 --- Notificación en paralelo, con los espectadores al final}

\begin{tabla}{|L{2.6cm}|Y|}
\fichacab
\parte{Estilo}{Táctica de desempeño: priorizar el trabajo.}
\parte{Descripción}{Al aceptar la jugada, el servicio notifica primero y a la vez a los dos participantes, y después a los espectadores, que CU-34 ya admite que reciban el estado con retraso. Atiende la parte de ESC-01 que exige que el oponente vea la jugada en el mismo plazo que quien la hizo. Se elige frente a notificar en secuencia a todos los destinatarios, donde el último espera a que terminen los anteriores y una partida con público retrasaría a sus propios jugadores.}
\parte{Efectos}{El número de espectadores deja de afectar al tiempo de la partida. Obliga a que el envío no bloquee al objeto activo de la partida (CD-05) mientras escribe en sockets lentos. Encaja con el retraso previsto que CU-21 menciona en su paso 20.}
\parte{Trazabilidad}{DRV-03, DRV-01. CD-05. ESC-01. CU-21 (paso 20), CU-20, CU-34. REQ-02. CRN-02. STK-01, STK-14.}
\end{tabla}

% ------------------------------------------------------------
\Needspace*{0.3\textheight}
\subsubsection*{CD-27 --- Medición por tramo incorporada al servidor}

\begin{tabla}{|L{2.6cm}|Y|}
\fichacab
\parte{Estilo}{Táctica de desempeño: supervisión del recurso.}
\parte{Descripción}{Cada jugada deja registrado cuánto tardó en cada tramo del camino, y el servidor expone esas cifras como distribución, con su mediana y su máximo. Atiende ESC-01 porque sin medición no hay forma de saber si el reparto de CD-23 se cumple, ni de encontrar el tramo culpable cuando no se cumple. Se elige frente a medir solo en pruebas de laboratorio, que no dicen nada del comportamiento con 20 partidas reales.}
\parte{Efectos}{La medida de ESC-01 se convierte en una prueba repetible del incremento semanal (CON-14). El registro por tramo no puede costar tanto como lo que mide, así que se acumula en memoria y se vuelca fuera del camino, apoyándose en CD-10. Sirve además de evidencia para STK-16.}
\parte{Trazabilidad}{DRV-03, DRV-08. CD-10, CD-23. ESC-01, ESC-17. CU-20, CU-21. REQ-02, CON-14. CRN-02, CRN-08. STK-01, STK-16, STK-14.}
\end{tabla}
